% =========================================================
% Basic Theorems of Rings
% =========================================================
\paragraph{Basic Theorems}

\begin{remark}[What Needs Proving]
The ring axioms say nothing explicitly about how multiplication
interacts with the additive identity $0$ or with additive inverses.
These must be derived. The key tool in both proofs is distributivity
--- the bridge between the two operations.

Notice that these proofs make no assumption about which ring you are in.
They work in $\mathbb{Z}$, in $M_n(\mathbb{R})$, in $\mathbb{Z}[x]$,
in any ring simultaneously.
\end{remark}

\begin{proposition}[Multiplication by Zero]
\label{prop:ring-mult-zero}
Let $R$ be a ring. For all $a \in R$,
\[
a \cdot 0 = 0 \cdot a = 0.
\]
\hyperref[prf:ring-mult-zero]{Go to proof.}
\end{proposition}

\begin{remark*}
\FlashQ{State: multiplication by zero in a ring.}
\FlashA{In any ring $R$: $a\cdot 0=0\cdot a=0$ for all $a\in R$.}
\FlashHint{Distributivity: $a\cdot 0=a(0+0)=a\cdot 0+a\cdot 0$; cancel by additive inverse.}
\end{remark*}



\begin{remark}
The proof uses distributivity to produce the equation
$a \cdot 0 + a \cdot 0 = a \cdot 0$, then applies additive
cancellation to conclude $a \cdot 0 = 0$.

This is not circular: $0$ on the left side is the additive identity
of the ring; the $0$ on the right is the same element being derived
as a consequence of cancellation. The proof works because additive
cancellation was already proved for all abelian groups
(Proposition~\ref{prop:group-cancellation}).
\end{remark}

\begin{proposition}[Multiplication by Additive Inverse]
\label{prop:ring-mult-neg}
Let $R$ be a ring. For all $a, b \in R$,
\[
a \cdot (-b) = -(a \cdot b)
\quad\text{and}\quad
(-a) \cdot b = -(a \cdot b).
\]
In any ring with unity, $(-1) \cdot a = -a$.
\hyperref[prf:ring-mult-neg]{Go to proof.}
\end{proposition}

\begin{remark*}
\FlashQ{In a ring, what is $a \cdot (-b)$?}
\FlashA{$a(-b) = -(ab)$ and $(-a)b = -(ab)$; in a ring with unity $(-1)a = -a$.}
\FlashHint{Show $a(-b)$ satisfies the defining property of $-(ab)$, then invoke uniqueness of additive inverse.}
\end{remark*}


\begin{remark}
The proof strategy is: show $a \cdot (-b)$ satisfies the defining
property of the additive inverse of $a \cdot b$, then invoke
uniqueness of additive inverses
(Proposition~\ref{prop:group-inverse-unique}).

This is the same strategy used in the vector space proof that
$(-1)v = -v$ --- because that proof is just this theorem applied
to the scalar field $\mathbb{F}$ acting on $V$.
\end{remark}
